\section{SonarQube -- Code Quality Analysis}
\label{sec:sonarqube}

\subsection{Giới thiệu}
SonarQube là nền tảng phân tích chất lượng code, phát hiện bugs, code smells, security vulnerabilities và đo lường coverage.

Nhóm sử dụng SonarQube \textbf{self-hosted} trên cùng AWS EC2 instance với Jenkins, truy cập tại \texttt{http://3.27.92.213:9000}.

\subsection{Cấu hình Server}
\begin{itemize}[noitemsep]
    \item \textbf{SonarQube URL}: \texttt{http://3.27.92.213:9000}
    \item \textbf{Authentication}: Token-based, lưu trong Jenkins Credentials (\texttt{sonarqube-token})
    \item \textbf{Project Key}: \texttt{nashtech-garage\_yas-yas-parent} (auto-detected từ pom.xml)
\end{itemize}

% [Ảnh: Screenshot cấu hình Jenkins System - SonarQube servers]
% [Ảnh: Screenshot cấu hình Webhook trên SonarQube]
% [Ảnh: SonarQube Dashboard - Overview]

\subsection{Tích hợp trong Pipeline}
\begin{lstlisting}[language=java, caption={SonarQube Analysis \& Quality Gate}]
stage('SonarQube - Analysis') {
    steps {
        withSonarQubeEnv('sonar-server') {
            withCredentials([string(credentialsId: 'sonarqube-token', variable: 'SONAR_TOKEN')]) {
                sh "mvn sonar:sonar -pl ${plArg} -am -Dsonar.host.url=..."
            }
        }
    }
}
stage('SonarQube - Quality Gate') {
    steps {
        withSonarQubeEnv('sonar-server') {
            timeout(time: 15, unit: 'MINUTES') {
                waitForQualityGate abortPipeline: false
            }
        }
    }
}
\end{lstlisting}

\subsection{Thiết kế và Quality Gate}
\begin{itemize}[noitemsep]
    \item Chỉ phân tích \textbf{service modules} thay đổi (exclude common-library)
    \item Sử dụng \texttt{withSonarQubeEnv} để Jenkins tự động nhận biết cấu hình server.
    \item Stage \textbf{Quality Gate} sử dụng lệnh \texttt{waitForQualityGate} để treo Pipeline cho đến khi SonarQube phân tích xong và gửi webhook kết quả về lại Jenkins.
\end{itemize}

\subsection{Kết quả}
\begin{itemize}[noitemsep]
    \item Console: \texttt{[INFO] ANALYSIS SUCCESSFUL}
    \item SonarQube task URL được in ra để truy cập dashboard
    \item Dashboard hiển thị: bugs, vulnerabilities, code smells, coverage percentage, duplications
\end{itemize}

% [Ảnh: SonarQube Dashboard - Project overview (bugs, coverage, smells)]
% [Ảnh: Console output "[INFO] ANALYSIS SUCCESSFUL"]
